% !TeX program = xelatex
% UTF-8 transcription of all seven source photographs.
% Wording and medical claims are preserved from the notes.
% pic01.jpg = 20261019_081315420.jpg (page 1)
% pic02.jpg = 20261019_081323300.jpg (page 2)
% pic03.jpg = 20261019_081328189.jpg (page 3)
% pic04.jpg = 20261019_081334219.jpg (page 4)
% pic05.jpg = 20261019_081349154.jpg (page 5)
% pic06.jpg = 20261019_081353923.jpg (page 6)
% pic07.jpg = 20261019_081403758.jpg (page 7)
% Compile directly with XeLaTeX, or input into a document using xepersian.
\ifdefined\settextfont
  \def\PEStandalone{0}
\else
  \def\PEStandalone{1}
  \documentclass[11pt,a4paper]{article}
  \usepackage[margin=1.8cm]{geometry}
  \usepackage{amsmath,amssymb,enumitem}
  \usepackage{xepersian}
  \IfFontExistsTF{Amiri}{\settextfont{Amiri}}{\settextfont{FreeSerif}}
  \setlatintextfont{FreeSerif}
  \setlength{\parindent}{0pt}
  \setlist{itemsep=0.35em,topsep=0.4em}
\fi
\providecommand{\en}[1]{\lr{#1}}
\newcommand{\PEarrow}{\en{$\longleftarrow$}}
\newcommand{\PEheader}[2][Year]{%
  {\footnotesize\en{#2\qquad Subject: 8\qquad
  #1: \_\_\_\quad Month: \_\_\_\quad Date: \_\_\_}\par}}
\newcommand{\PEfooter}[2]{%
  \par\medskip{\footnotesize\en{#1}\hfill #2\par}}
\ifnum\PEStandalone=1 \begin{document}\fi

% Source page 1
\PEheader{}
\section*{\en{Pleural Effusion}}
\begin{enumerate}[label=\arabic*.,leftmargin=*]
\item در فضای پلور به صورت نرمال \en{2--5 ml} مایع پلورال وجود دارد.
\item مایع پلور در پلور پریتال تولید می‌شود، ولی توسط هر دو پلور جذب می‌شود.
\item \en{CXR} معمولاً \textbf{اولین قدم} در تشخیص پلورال افیوژن است.
\item \textbf{سونو و \en{CT}} می‌توانند مقادیر اندک پلورال افیوژن را نشان دهند.
\item \en{CT} بهترین تصویربرداری برای بررسی بیماری ریهٔ عامل افیوژن است.
  \begin{itemize}
  \item همچنین بهترین تصویربرداری در شرایطی که یک همی‌توراکس به علت افیوژن کاملاً اوپاسیفیه شده است.
  \end{itemize}
\item زمانی افیوژن ایجاد می‌شود که تولید مایع پلورال از جذب آن بیشتر شود.
\item علل پلورال افیوژن
  \begin{itemize}
  \item افزایش فشار هیدرواستاتیک
  \item کاهش فشار انکوتیک: هایپوپروتئینمی
  \item افزایش نفوذپذیری مویرگ‌ها
    \begin{itemize}
    \item سموم \PEarrow\ \en{Toxic disruption of capillary membrane}
    \item واکنش‌های ازدیاد حساسیت
    \end{itemize}
  \end{itemize}
\end{enumerate}
\PEfooter{Batus}{۱}

\clearpage
% Source page 2
\PEheader{}
\begin{enumerate}[label=\arabic*.,leftmargin=*]
\setcounter{enumi}{7}
\item علل پلورال افیوژن
  \begin{itemize}
  \item افزایش تولید
    \begin{itemize}
    \item \en{CHF}
    \item هایپوپروتئینمی
    \item افیوژن پاراپنومونیک
    \item ازدیاد حساسیت
    \end{itemize}
  \item کاهش بازجذب
    \begin{itemize}
    \item انسداد لنفاوی ناشی از تومور
    \item افزایش فشار ورید مرکزی
    \item کاهش فشار اینتراپلورال \PEarrow\ \textbf{آتلکتازی یا انسداد برونش}
    \end{itemize}
  \item انتقال از حفرهٔ پریتوئن \PEarrow\ آسیت
  \end{itemize}
\item انواع پلورال افیوژن
  \begin{itemize}
  \item اگزوداتیو
    \begin{itemize}
    \item ناشی از یک پروسهٔ التهابی
    \item شایع‌ترین علت، بدخیمی است.
    \item آمپیم
    \item هموتوراکس
    \item شیلوتوراکس
    \end{itemize}
  \item ترانسوداتیو
    \begin{itemize}
    \item \en{CHF} (عموماً نارسایی قلب چپ): شایع‌ترین علت
    \item هایپوآلبومینمی
    \item سیروز
    \item سندرم نفروتیک
    \end{itemize}
  \end{itemize}
\end{enumerate}
\PEfooter{Batus}{۲}

\clearpage
% Source page 3
\PEheader{}
\begin{enumerate}[label=\arabic*.,leftmargin=*]
\setcounter{enumi}{9}
\item پلورال افیوژن دوطرفه
  \begin{itemize}
  \item \en{CHF} \PEarrow\ دوطرف، به یک میزان و یا راست کمی بیشتر
  \item لوپوس \PEarrow\ معمولاً دوطرفه
    \begin{itemize}
    \item اگر یک‌طرفه باشد معمولاً سمت چپ
    \end{itemize}
  \item \textbf{اگر یک طرف خیلی بیشتر بود} \PEarrow\ بدخیمی، عفونت
  \end{itemize}
% The source numbering jumps from 10 to 21; preserve it verbatim.
\setcounter{enumi}{20}
\item پلورال افیوژن یک‌طرفه
  \begin{itemize}
  \item عفونت‌ها، \en{TB}
  \item تروما
  \item بیماری ترومبوآمبولیک ریوی
  \end{itemize}
\item بیماری‌هایی که معمولاً \en{PE} سمت چپ می‌دهند
  \begin{itemize}
  \item پانکراتیت
  \item انسداد دیستال توراسیک داکت
  \item \en{Dressler's Syndrome}
    \begin{latin}
    \noindent $\longrightarrow$ Postpericardiotomy/postmyocardial infarction syndrome
    \end{latin}
    \begin{itemize}
    \item \en{2--3} هفته بعد از میوکاردیت ترانسمورال
    \item درمان با آسپرین یا استروئید \en{high-dose}
    \end{itemize}
  \end{itemize}
\item بیماری‌هایی که معمولاً \en{PE} سمت راست می‌دهند:
  \begin{itemize}
  \item بیماری‌های شکمی مرتبط با کبد و تخمدان
  \item برخی تومورهای تخمدان \PEarrow\ آسیت و \en{PE} راست
    \begin{itemize}
    \item \en{Meig's Syndrome}
    \end{itemize}
  \item \en{RA}: \en{PE} که سال‌ها بدون تغییر می‌ماند.
  \item انسداد پروگزیمال توراسیک داکت
  \end{itemize}
\end{enumerate}
\PEfooter{Batus}{۳}

\clearpage
% Source page 4
\PEheader{}
\begin{enumerate}[label=\arabic*.,leftmargin=*]
\setcounter{enumi}{23}
\item افیوژن ساب‌پولمونیک
  \begin{itemize}
  \item تقریباً همهٔ افیوژن‌ها ابتدا در فضای ساب‌پولمونیک در زیر ریه تجمع پیدا می‌کنند.
  \item اگر فقط افیوژن ساب‌پولمونیک داشته باشیم، تشخیص در \en{CXR} دشوار است.
  \item اغلب افیوژن‌های ساب‌پولمونیک آزاد هستند و با تغییر پوزیشن بیمار جابجا می‌شوند.
  \end{itemize}
\item بلانت شدن سولکوس کاستوفرنیک خلفی
  \begin{itemize}
  \item عموماً اولین یافتهٔ پلورال افیوژن
  \item در گرافی لترال: به \en{75 cc} مایع پلورال نیاز است.
  \item در گرافی \en{PA}: به \en{300 cc} مایع پلورال
  \end{itemize}
\item ضخیم شدن پلور ناشی از فیبروز می‌تواند باعث بلانت شدن سولکوس شود.
  \begin{itemize}
  \item \textbf{\en{Ski-slope appearance}}
  \item با تغییر پوزیشن تغییر نمی‌کند.
  \end{itemize}
\item \textbf{\en{Meniscus Sign}}
  \begin{itemize}
  \item مایع پلورال در طرف قفسهٔ سینه بالاتر قرار می‌گیرد.
    \begin{latin}
    \noindent $\longrightarrow$ higher on the side, lower in the middle
    \end{latin}
  \end{itemize}
\end{enumerate}
\PEfooter{Batus}{۴}

\clearpage
% Source page 5: the printed year field reads "Veor".
\PEheader[Veor]{Note book}
\begin{enumerate}[label=\arabic*.,leftmargin=*]
\setcounter{enumi}{27}
\item در گرافی سوپاین، \en{PE} در پشت ریه جمع شده و یک نمای هموژن \en{haze} در تمام ریه ایجاد می‌کند.
\item گرافی دکوبیتوس \PEarrow\ بفهمیم مایع افیوژن آزادانه حرکت دارد یا لوکوله است.
  \begin{itemize}
  \item برای درناژ مهم است.
  \end{itemize}
\item گرافی دکوبیتوس می‌تواند \en{PE} را در حد \en{15--20 ml} تشخیص بدهد.
\item اگر در گرافی \en{PA} تمام ریه اوپاسیفیه باشد گرافی دکوبیتوس کمکی نمی‌کند.
  \begin{itemize}
  \item \en{CT} می‌گیریم.
  \end{itemize}
\item \en{Opacified Hemithorax}
  \begin{itemize}
  \item یعنی حداقل \en{2 L} مایع در همی‌توراکس است.
  \item باعث \en{passive collapse} ریه (آتلکتازی) می‌شود.
  \item اندیکاسیون \en{CT}
  \end{itemize}
\item افیوژن \en{large} می‌تواند قلب و تراشه را جابجا کند.
\item افیوژن لوکوله
  \begin{itemize}
  \item اگر یک عفونت قبلی ارگانیزه‌شده باشد.
  \item زمانی شک به لوکولیشن می‌کنیم که افیوژن:
    \begin{latin}
    \noindent $\longrightarrow$ Unusual shape or location
    \end{latin}
  \end{itemize}
\end{enumerate}
\PEfooter{BAHAR DANESH ASL}{۵}

\clearpage
% Source page 6
\PEheader[Veor]{Note book}
\begin{enumerate}[label=\arabic*.,leftmargin=*]
\setcounter{enumi}{34}
\item سودوتومور (\en{Vanishing Tumors}) (فیشرال سودوتومور)
  \begin{latin}
  \begin{itemize}
  \item Sharply marginated collections of PE
  \end{itemize}
  \end{latin}
  \begin{itemize}
  \item تقریباً همیشه در بیماران \en{CHF}
    \begin{itemize}
    \item همیشه ترانسوداتیو
    \end{itemize}
  \end{itemize}
  \begin{latin}
  \begin{itemize}
  \item They are lenticular in shape
  \item 75\% in minor fissure
  \item Pointed ends on each side
    \begin{itemize}
    \item Shape of a lemon
    \end{itemize}
  \end{itemize}
  \end{latin}
  \begin{itemize}
  \item با تغییر پوزیشن حرکت نمی‌کند.
  \item با درمان \en{CHF} رفع می‌شوند ولی اگر \en{CHF} برگردد دوباره در همان محل عود می‌کنند.
  \end{itemize}
\item لامینار افیوژن
  \begin{latin}
  \begin{itemize}
  \item Thin, band like, along the lateral chest wall especially near the costophrenic angle.
  \end{itemize}
  \end{latin}
  \begin{itemize}
  \item \textbf{تقریباً همواره ناشی از افزایش فشار دهلیز چپ است.}
    \begin{itemize}
    \item \en{CHF}
    \item \en{Lymphangitic spread of Malignancy}
    \end{itemize}
  \end{itemize}
  \begin{latin}
  \noindent $\longrightarrow$ Usually not free fluid
  \end{latin}
\end{enumerate}
\PEfooter{BAHAR DANESH ASL}{۶}

\clearpage
% Source page 7
\PEheader[Veor]{Note book}
\begin{enumerate}[label=\arabic*.,leftmargin=*]
\setcounter{enumi}{36}
\item هیدروپنوموتوراکس
  \begin{itemize}
  \item وجود همزمان مایع و هوا
  \item علل
    \begin{itemize}
    \item تروما
    \item جراحی
    \item توراسنتز اخیر برای تخلیهٔ \en{PE}
    \end{itemize}
  \item فیستول برونکوپلورال نیز می‌تواند عامل آن باشد.
    \begin{itemize}
    \item ناشایع است.
    \item علل
      \begin{itemize}
      \item تومور
      \item جراحی
      \item عفونت
      \end{itemize}
    \end{itemize}
  \item برای افتراق از یک آبسهٔ بزرگ \PEarrow\ حتماً \en{CT}
  \item سطح مایع -- هوا
  \end{itemize}
\end{enumerate}
\PEfooter{BAHAR DANESH ASL}{۷}

\ifnum\PEStandalone=1 \expandafter\enddocument\fi
